\documentclass[conference]{IEEEtran}
\IEEEoverridecommandlockouts

\usepackage{cite}
\usepackage{amsmath,amssymb,amsfonts}
\usepackage{algorithmic}
\usepackage{algorithm}
\usepackage{graphicx}
\usepackage{textcomp}
\usepackage{xcolor}
\usepackage{booktabs}
\usepackage{array}

\begin{document}

\title{Constant Bandwidth Server: Technical Briefing\\
{\footnotesize Milestone 1 --- Real-Time Systems Seminar, HSHL}}

\author{\IEEEauthorblockN{Soaib Ferdous}
\IEEEauthorblockA{\textit{Department of Computer Science} \\
\textit{Hochschule Hamm-Lippstadt}\\
Student ID: 1232368 \\
Group B, Team 6}}

\maketitle

\begin{abstract}
This document is the Milestone 1 technical briefing for the seminar
topic \emph{Constant Bandwidth Server} (CBS). It covers the problem
setting that motivates CBS, the formal definition and budget/deadline
management mechanism, the scheduling algorithm, assumptions, runtime
overhead, and open questions identified during study of the primary
reference (Buttazzo, \textit{Hard Real-Time Computing Systems}, 4th ed.).
\end{abstract}

% ============================================================
\section{Problem Statement and Motivation}
% ============================================================

\begin{itemize}
  \item Real-time systems typically mix \emph{periodic} tasks (hard
        deadlines, known periods) with \emph{aperiodic} tasks (irregular
        arrivals, often soft deadlines).
  \item Under fixed-priority schedulers (RM/DM), naive background
        execution of aperiodic tasks yields poor response times;
        foreground execution can cause periodic tasks to miss deadlines.
  \item Earlier server approaches attempt a middle ground:
        \emph{Polling Server} (PS) wastes budget if no aperiodic job is
        waiting; \emph{Deferrable Server} (DS) preserves budget but can
        violate the RM schedulability bound because it behaves like a
        task with higher effective utilization at replenishment points.
  \item \textbf{Core problem:} How can aperiodic tasks be served with
        low response time \emph{and} with a provable guarantee that
        periodic task schedulability is not affected?
  \item CBS answers this under EDF by ensuring the server never
        increases processor demand beyond what is implied by its
        assigned bandwidth $U_s = Q_s / T_s$.
\end{itemize}

% ============================================================
\section{Core Technical Contribution}
% ============================================================

\begin{itemize}
  \item CBS is a dynamic-priority server designed to work with
        \emph{Earliest Deadline First} (EDF) scheduling.
  \item Key property: \textbf{temporal isolation} --- a misbehaving or
        overloaded aperiodic task cannot steal CPU time from periodic
        tasks beyond what the server bandwidth allows.
  \item Mechanism: the server maintains a \emph{current deadline}
        $d_s$ and a \emph{current budget} $c_s$. An aperiodic job runs
        only when $d_s$ is the earliest deadline in the system. When
        budget is exhausted the deadline is postponed and budget is
        recharged --- the task continues but at lower effective priority.
  \item This distinguishes CBS from DS: DS can retain budget across a
        period boundary, increasing instantaneous demand. CBS prevents
        this via its replenishment rule (Section~\ref{sec:algo}).
\end{itemize}

% ============================================================
\section{Formal Definition}
% ============================================================

A CBS is defined by two parameters:
\begin{itemize}
  \item $Q_s$ --- maximum server budget (execution time per period)
  \item $T_s$ --- server period
\end{itemize}

\noindent\textbf{Server bandwidth:}
\begin{equation}
  U_s = \frac{Q_s}{T_s}
  \label{eq:bandwidth}
\end{equation}

\noindent\textbf{Schedulability condition} (CBS + $n$ periodic EDF tasks):
\begin{equation}
  U_s + \sum_{i=1}^{n} \frac{C_i}{T_i} \leq 1
  \label{eq:sched}
\end{equation}
where $C_i$ and $T_i$ are the WCET and period of periodic task
$\tau_i$. CBS is treated identically to a periodic task with
utilization $U_s$ in the EDF schedulability test --- this is its
central analytical advantage over DS.

\noindent\textbf{Runtime server state variables:}
\begin{itemize}
  \item $c_s$: remaining budget, $0 \leq c_s \leq Q_s$
  \item $d_s$: current server deadline
\end{itemize}

% ============================================================
\section{Algorithm and Budget/Deadline Rules}
\label{sec:algo}
% ============================================================

\subsection{Initialization}
At $t = 0$: set $c_s \leftarrow Q_s$ and $d_s \leftarrow T_s$.

\subsection{Job Arrival Rule}
When aperiodic job $J$ arrives at time $t$:
\begin{itemize}
  \item \textbf{If} server is idle and $c_s \geq (d_s - t) \cdot U_s$:
        assign current $d_s$ to $J$ (budget is still fresh enough).
  \item \textbf{Else:} set $d_s \leftarrow d_s + T_s$,
        $c_s \leftarrow Q_s$; assign new $d_s$ to $J$.
  \item \textbf{If} server is busy: $J$ is queued (FIFO).
\end{itemize}

The condition $c_s \geq (d_s - t) \cdot U_s$ prevents using stale
budget that would cause demand to exceed $U_s$ within the current
period.

\subsection{Budget Exhaustion Rule}
When $c_s = 0$ during execution at time $t$:
set $d_s \leftarrow d_s + T_s$ and $c_s \leftarrow Q_s$.
The job continues with the new (later) $d_s$, reducing its EDF
priority automatically.

\begin{algorithm}
\caption{CBS scheduling decisions (simplified)}
\begin{algorithmic}[1]
\STATE \textbf{On job arrival} at time $t$:
\IF{server idle \AND $c_s \geq (d_s - t) \cdot U_s$}
  \STATE assign current $d_s$ to job
\ELSE
  \STATE $d_s \leftarrow d_s + T_s$;\ $c_s \leftarrow Q_s$
  \STATE assign new $d_s$ to job
\ENDIF
\STATE \textbf{On budget exhaustion} ($c_s = 0$):
\STATE $d_s \leftarrow d_s + T_s$;\ $c_s \leftarrow Q_s$
\STATE job continues with updated $d_s$
\end{algorithmic}
\end{algorithm}

% ============================================================
\section{Assumptions and Scope Limitations}
% ============================================================

\begin{itemize}
  \item \textbf{Scheduler:} EDF required. CBS does not apply to
        fixed-priority schedulers (RM/DM). The Sporadic Server is the
        fixed-priority equivalent.
  \item \textbf{Processor model:} Single processor in the original
        Buttazzo formulation. Multi-core extension is non-trivial.
  \item \textbf{Preemption:} Fully preemptive model assumed throughout.
  \item \textbf{Task model:} Periodic tasks with implicit deadlines
        ($D_i = T_i$) and known $C_i$, $T_i$.
  \item \textbf{Aperiodic guarantees:} Bandwidth only --- no hard
        per-job response time guarantee is provided.
  \item \textbf{No resource sharing:} Priority inversion and shared
        resources are outside the scope of basic CBS.
\end{itemize}

% ============================================================
\section{Runtime and Overhead}
% ============================================================

\begin{itemize}
  \item \textbf{Server state update:} $O(1)$ per arrival event and
        per budget exhaustion event (comparison + addition only).
  \item \textbf{EDF queue:} $O(\log n)$ insert/update, shared with
        all tasks in the system.
  \item \textbf{Timer:} One additional timer interrupt per budget
        period when the server is active (to detect $c_s = 0$).
  \item \textbf{Memory:} Constant per server instance: stores $Q_s$,
        $T_s$, $c_s$, $d_s$, and a job queue pointer.
  \item \textbf{Comparison to DS:} CBS adds one conditional check at
        job arrival not present in DS. Overhead is negligible relative
        to the scheduling correctness gained.
\end{itemize}

% ============================================================
\section{Relevance for Embedded Systems}
% ============================================================

\begin{itemize}
  \item Embedded control systems frequently combine periodic control
        tasks (hard RT) with event-driven handlers (soft RT, e.g.
        sensor interrupts, communication events).
  \item CBS allows a single EDF scheduler to handle both task classes
        with a formal bandwidth separation guarantee.
  \item The $O(1)$ server-state overhead is compatible with
        resource-constrained processors, though EDF itself requires
        more runtime support than RM.
  \item Temporal isolation is critical in safety-relevant systems
        where one task class must not interfere with another under
        any load condition.
\end{itemize}

% ============================================================
\section{Open Questions for Milestone 2}
% ============================================================

\begin{enumerate}
  \item \textbf{Hard CBS variant:} Buttazzo describes a stricter
        replenishment rule. How does it differ and does it affect the
        schedulability bound?
  \item \textbf{Response time bound:} Is there a closed-form worst-case
        aperiodic response time under CBS, analogous to DS under RM?
  \item \textbf{CBS vs. TBS:} Total Bandwidth Server also works under
        EDF with a simpler rule. What exactly does CBS provide that TBS
        does not? (To be analyzed in M2.)
  \item \textbf{Multiple CBS instances:} Can multiple CBS servers
        coexist on one processor? How does the schedulability condition
        generalize?
  \item \textbf{UPPAAL modeling:} How should the budget replenishment
        transition be modeled as a timed automaton? Which state
        variables map to clock constraints?
\end{enumerate}

% ============================================================
\section*{References}
% ============================================================

\begin{thebibliography}{9}

\bibitem{buttazzo}
G.~C.~Buttazzo,
\textit{Hard Real-Time Computing Systems: Predictable Scheduling
Algorithms and Applications}, 4th~ed.
Springer, 2024.
[Primary reference. CBS is covered in the chapter on aperiodic task
servers. Page numbers to be verified upon book access.]

\bibitem{abeni1998}
L.~Abeni and G.~Buttazzo,
``Integrating Multimedia Applications in Hard Real-Time Systems,''
in \textit{Proc. 19th IEEE Real-Time Systems Symposium
(Cat. No. 98CB36279)}, IEEE, 1998, pp.~4--13.
[Original CBS paper introducing the formal definition, budget and
deadline rules, and schedulability proof under EDF.]

\bibitem{abeni2004}
L.~Abeni and G.~Buttazzo,
``Resource Reservation in Dynamic Real-Time Systems,''
\textit{Real-Time Systems}, vol.~27, no.~2, pp.~123--167, 2004.
[Extended analysis of CBS including response time results and
comparison with TBS.]

\bibitem{ai_usage_m1}
S.~Ferdous,
\textit{AI Usage Protocol, Milestone 1: Constant Bandwidth Server},
HSHL, 2026.
[Protocol version 0.2; tool: ChatGPT GPT-5.5; 6 documented entries;
JSON and .bib files submitted alongside this document.]

\end{thebibliography}

\end{document}
